\begin{frame}
	\frametitle{Como treinar uma função que não tem gradiente?}
	\only<1>{
		\framesubtitle{O mesmo problema estrutural do BitNet}
		\par A função de disparo do LIF se aproxima de uma \textbf{função em degrau} (Heaviside): dispara ou não dispara, sem meio-termo. Sua derivada é zero em quase todo ponto --- a descida de gradiente padrão não tem o que propagar \cite{kasabov2019time}.
		\par Quatro famílias de resposta, historicamente:
		\begin{itemize}
			\item \textbf{STDP}: regra local --- se o neurônio pré-sináptico dispara \textbf{antes} do pós-sináptico, a conexão se fortalece; se depois, se enfraquece. Biologicamente plausível, não supervisionada.
			\item \textbf{Gradiente substituto + BPTT}: usa a função em degrau real no \textit{forward}, e uma aproximação suave e diferenciável só no \textit{backward} --- a mesma lógica do STE do BitNet.
			\item \textbf{Algoritmos evolutivos}: seleção dos hiperparâmetros mais aptos ao longo de gerações.
			\item \textbf{Reservatório / computação dinâmica}: treina só a camada de saída, mantendo pesos recorrentes fixos (\textit{echo state networks}, \textit{liquid state machines}).
		\end{itemize}
	}
	\only<2>{
		\framesubtitle{Alternativa: converter em vez de treinar}
		\begin{itemize}
			\item Treina uma \textbf{RNA convencional}, com técnicas padrão de otimização.
			\item \textbf{Converte} os pesos treinados para uma RNP, balanceando pesos e limiares de disparo \cite{diehl2015fast, rueckauer2017conversion}.
			\item Evita o problema do gradiente por completo --- mas herda limitações: geralmente precisa de \textbf{mais passos de tempo} para aproximar a acurácia da RNA original, e nem toda arquitetura converte bem.
		\end{itemize}
	}
	% [SOFTWARE] SNN -> Surrogate gradient | checkpoints "A derivada
	% real" -> "O gradiente substituto" (sweep sigmoide+gradiente)
	\abrirNoSoftware{SNN \textrightarrow{} Surrogate gradient}{snn.surrogate}
\end{frame}
